En un modelo determinista de circulación, las cantidades \(x\) e \(y\) en dos estados A y B se transfieren en ambos sentidos con la misma tasa \(k=1\,\mathrm{min}^{-1}\):
\[
x'=-kx+ky,\qquad y'=kx-ky.
\]
Se comparan dos preparaciones sintéticas, P con \((x(0),y(0))=(8,2)\) y Q con \((2,8)\). Un monitor informa de la cantidad total \(x+y\). Se puede guardar además uno de estos registros: un contador que reúne los traslados en ambos sentidos, o un contador que conserva sólo los traslados de A a B.

Los contadores empiezan en cero en cada ensayo y acumulan las integrales de las tasas de traslado; no son recuentos aleatorios. Un elemento puede trasladarse y contarse varias veces. Para una duración \(T>0\), determine qué información permiten distinguir el monitor total y cada contador, y decida qué registro adicional conserva la diferencia entre P y Q.

Compare también el límite de la diferencia entre sus cantidades instantáneas con el límite de la diferencia entre sus contadores direccionales cuando \(T\) aumenta. Distinga una diferencia histórica finita del valor de un contador que sigue acumulando eventos.
